\documentclass{ximera}

\begin{document}
	\author{Stitz-Zeager}
	\xmtitle{Exercises for Absolute Value Functions}{}

\mfpicnumber{1} \opengraphsfile{ExercisesforAbsoluteValueFunctions} % mfpic settings added 

\begin{question}
In Exercises \ref{graphbasicabsvalexerfirst} - \ref{graphbasicabsvalexerlast}, graph the function using Theorem \ref{linearabsvaluegraphs}. Find the axis intercepts of each graph, if any exist.  From the graph, determine the domain and range of each function, the maximum and minimum of each function, if they exist, and list the intervals on which the function is increasing, decreasing or constant.

\begin{problem}\label{graphbasicabsvalexerfirst}
$f(x) = |x + 4|$

\begin{solution}
Use the Desmos graph below with the settings $v(x)=|x|,\,h=-4,\,k=0,\,a=1$

\begin{center} 
\desmos{hmwcgnrph3}{800}{600} 
\end{center}

$x$-intercept $(-4, 0)$ \\ $y$-intercept $(0, 4)$ \\ Domain $(-\infty, \infty)$ \\ Range $[0, \infty)$ \\ Decreasing on $(-\infty, -4]$ \\ Increasing on $[-4, \infty)$ \\ Minimum is $0$ at $(-4,0)$ \\ No maximum\\

\end{solution}

\end{problem}

\begin{problem}
$f(x) = |x| + 4$

\begin{solution}
Use the Desmos graph below with the settings $v(x)=|x|,\,h=0,\,k=4,\,a=1$

\begin{center} 
\desmos{hmwcgnrph3}{800}{600} 
\end{center}

No $x$-intercepts \\ $y$-intercept $(0, 4)$ \\ Domain $(-\infty, \infty)$ \\ Range $[4, \infty)$ \\ Decreasing on $(-\infty, 0]$ \\ Increasing on $[0, \infty)$ \\  Minimum is $4$ at $(0,4)$  \\ No maximum \\

\end{solution}
\end{problem}

\begin{problem}
$f(x) = |4x|$

\begin{solution}
Note $f(x) = |4x| = 4|x|$.

Use the Desmos graph below with the settings $v(x)=|x|,\,h=0,\,k=0,\,a=4$

\begin{center} 
\desmos{hmwcgnrph3}{800}{600} 
\end{center}

$x$-intercept $(0, 0)$ \\ $y$-intercept $(0, 0)$ \\ Domain $(-\infty, \infty)$ \\ Range $[0, \infty)$ \\ Decreasing on $(-\infty, 0]$ \\ Increasing on $[0, \infty)$ \\ Minimum is $0$ at $(0,0)$  \\ No maximum \\
\end{solution}
\end{problem}

\begin{problem}
$g(t) = -3|t|$ 

\begin{solution}
Use the Desmos graph below with the settings $v(x)=|x|,\,h=0,\,k=0,\,a=-3$

\begin{center} 
\desmos{hmwcgnrph3}{800}{600} 
\end{center}

$t$-intercept $(0, 0)$ \\ $y$-intercept $(0, 0)$ \\ Domain $(-\infty, \infty)$ \\ Range $(-\infty, 0]$ \\ Increasing on $(-\infty, 0]$ \\ Decreasing on $[0, \infty)$ \\ Maximum is $0$ at $(0, 0)$ \\ No minimum \\
\end{solution}
\end{problem}

\begin{problem}
$g(t) = 3|t + 4| - 4$

\begin{solution}
Use the Desmos graph below with the settings $v(x)=|x|,\,h=-4,\,k=-4,\,a=3$

\begin{center} 
\desmos{hmwcgnrph3}{800}{600} 
\end{center}

$t$-intercepts $\left(-\frac{16}{3}, 0\right)$, $\left(-\frac{8}{3}, 0\right)$ \\ $y$-intercept $(0, 8)$ \\ Domain $(-\infty, \infty)$ \\ Range $[-4, \infty)$ \\ Decreasing on $(-\infty, -4]$ \\ Increasing on $[-4, \infty)$ \\ Minimum is $-4$ at $(-4,-4)$ \\ No maximum \\ 
\end{solution}
\end{problem}
  
\begin{problem}\label{graphbasicabsvalexerlast}
$g(t) = \frac{1}{3}|2t - 1|$

\begin{solution}
Note $g(t) = \frac{1}{3}|2t - 1| = \frac{2}{3}|t - \frac{1}{2}|$

Use the Desmos graph below with the settings $v(x)=|x|,\,h=\frac{1}{2},\,k=0,\,a=\frac{2}{3}$

\begin{center} 
\desmos{hmwcgnrph3}{800}{600} 
\end{center}

$t$-intercepts $\left(\frac{1}{2}, 0\right)$ \\ $y$-intercept $\left(0, \frac{1}{3}\right)$ \\ Domain $(-\infty, \infty)$ \\ Range $[0, \infty)$ \\ Decreasing on $\left(-\infty, \frac{1}{2}\right]$ \\ Increasing on $\left[\frac{1}{2}, \infty\right)$ \\ Minimum is $0$ at $\left(\frac{1}{2},0\right)$ \\ No maximum \\ 
\end{solution}
\end{problem}
\end{question}

\begin{question}
In Exercises \ref{findformulaforabsgraphfirst} - \ref{findformulaabsgraphlast}, find a formula for each function below in the form $F(x) = a|x-h|+k$.

\begin{problem}\label{findformulaforabsgraphfirst}
\begin{tikzpicture}
\begin{axis}[
    width=3in, height=3in,
    xmin=-5, xmax=5, ymin=-5, ymax=5,
    axis lines=middle,
    axis line style={-{Latex[length=8pt,width=6pt]}},
    tick style={line width=1pt},
    major tick length=5pt,
    xlabel={$x$}, ylabel={$y$},
    xlabel style={anchor=west},
    ylabel style={anchor=south},
    xtick={-4,-3,-2,-1,0,1,2,3,4},
    ytick={-4,-3,-2,-1,0,1,2,3,4},
    tick label style={font=\tiny},
    label style={font=\scriptsize},
    samples=100,
]
\addplot[<->,line width=2.5pt, blue, domain=-4.5:2.5] {2*abs(x+1)-3};
\addplot[only marks,mark=*,mark size=4pt] coordinates {(0,-1)(-1,-3)};
\node at (0,-1) [right=4pt] {$(0,-1)$};
\node at (-1,-3) [below=8pt,left=8pt] {$(-1,-3)$};
\end{axis}
\node at (rel axis cs:0.5,0) [below] {The graph of $y=F(x)$.};
\end{tikzpicture}

\begin{solution}
$F(x) = 2|x+1|-3$
\end{solution}
\end{problem}

\begin{problem}
\begin{tikzpicture}
\begin{axis}[
    width=3in, height=3in,
    xmin=-5, xmax=5, ymin=-5, ymax=5,
    axis lines=middle,
    axis line style={-{Latex[length=8pt,width=6pt]}},
    tick style={line width=1pt},
    major tick length=5pt,
    xlabel={$x$}, ylabel={$y$},
    xlabel style={anchor=west},
    ylabel style={anchor=south},
    xtick={-4,-3,-2,-1,0,1,2,3,4},
    ytick={-4,-3,-2,-1,0,1,2,3,4},
    tick label style={font=\tiny},
    label style={font=\scriptsize},
    samples=100,
]
\addplot[<->,line width=2.5pt, blue, domain=-4.5:4.5] {abs(x-1.25)-2.75};
\addplot[only marks,mark=*,mark size=4pt] coordinates {(4,0)(-1.5,0)(0,-1.5)};
\node at (4,0) [above=4pt,left=4pt] {$(4,0)$};
\node at (-1.5,0) [above=4pt,left=8pt] {$(-1.5,0)$};
\node at (0,-1.5) [left=8pt] {$(0,-1.5)$};
\end{axis}
\node at (rel axis cs:0.5,0) [below] {The graph of $y=F(x)$.};
\end{tikzpicture}

\begin{solution}
$F(x) = |x-1.25|-2.75$
\end{solution}
\end{problem}

\begin{problem}
\begin{tikzpicture}
\begin{axis}[
    width=3in, height=3in,
    xmin=-5, xmax=5, ymin=-5, ymax=5,
    axis lines=middle,
    axis line style={-{Latex[length=8pt,width=6pt]}},
    tick style={line width=1pt},
    major tick length=5pt,
    xlabel={$x$}, ylabel={$y$},
    xlabel style={anchor=west},
    ylabel style={anchor=south},
    xtick={-4,-3,-2,-1,0,1,2,3,4},
    ytick={-4,-3,-2,-1,0,1,2,3,4},
    tick label style={font=\tiny},
    label style={font=\scriptsize},
    samples=100,
]
\addplot[<->,line width=2.5pt, blue, domain=-4.5:4.5] {-abs(x+1)+2};
\addplot[only marks,mark=*,mark size=4pt] coordinates {(-3,0)(1,0)(0,1)};
\node at (-3,0) [above=8pt,left=4pt] {$(-3,0)$};
\node at (1,0) [above=8pt,right=6pt] {$(1,0)$};
\node at (0,1) [above=10pt,right=6pt] {$(0,1)$};
\end{axis}
\node at (rel axis cs:0.5,0) [below] {The graph of $y=F(x)$.};
\end{tikzpicture}
\end{problem}

\begin{solution}
$F(x) = -|x+1|+2$
\end{solution}
\end{question}

\begin{problem}
\begin{tikzpicture}
\begin{axis}[
    width=3in, height=3in,
    xmin=-5, xmax=5, ymin=-5, ymax=5,
    axis lines=middle,
    axis line style={-{Latex[length=8pt,width=6pt]}},
    tick style={line width=1pt},
    major tick length=5pt,
    xlabel={$x$}, ylabel={$y$},
    xlabel style={anchor=west},
    ylabel style={anchor=south},
    xtick={-4,-3,-2,-1,0,1,2,3,4},
    ytick={-4,-3,-2,-1,0,1,2,3,4},
    tick label style={font=\tiny},
    label style={font=\scriptsize},
    samples=100,
]
\addplot[<->,line width=2.5pt, blue, domain=-4.5:4.5] {-1/2*abs(x-1)+3/2};
\addplot[only marks,mark=*,mark size=4pt] coordinates {(-2,0)(4,0)(0,1)(1,1.5)};
\node at (-2,0) [above=8pt,left=4pt] {$(-2,0)$};
\node at (4,0) [above=10pt] {$(4,0)$};
\node at (0,1) [above=4pt,left=8pt] {$(0,1)$};
\end{axis}
\node at (rel axis cs:0.5,0) [below] {The graph of $y=F(x)$.};
\end{tikzpicture}

\begin{solution}
$F(x) = -\frac{1}{2} |x+1|+\frac{3}{2}$
\end{solution}
\end{problem}

\begin{problem}\label{makeaveewithabsval}
Use\href{https://www.desmos.com/calculator/tzb1yozj6m}{this Desmos link} to graph the following pairs of functions on the same set of axes.  (The first one is done for you, and you can just modify f(x) to get the others.)

\begin{itemize}

\item  $f(x) = 2-x$ and $g(x) = | 2-x |$



\item  $f(x) = x^2-4$ and $g(x) = | x^2 -4 |$

\item  $f(x) = x^3$ and $g(x) = | x^3 |$

\item  $f(x) = \sqrt{x}-4 $ and $g(x) = | \sqrt{x} -4| $

\end{itemize}

Choose more functions $f(x)$ and graph $y = f(x)$ alongside $y = | f(x)|$ until you can explain how, in general, one would obtain the graph of $y = | f(x) |$ given the graph of $y = f(x)$.  How does your explanation tie in with  with Definition \ref{absolutevaluepiecewise}?

\begin{solution}
In each case, the graph of $g$ can be obtained from the graph of $f$ by reflecting the portion of the graph of $f$ which lies below the $x$-axis about the $x$-axis.  This meshes with Definition \ref{absolutevaluepiecewise} since what we are doing algebraically is making the negative $y$-values positive.
\end{solution}
\end{problem}

\begin{problem}
Explain why the function below cannot be written in the form $F(x) = a|x-h|+k$.  Write $F(x)$ as a piecewise-defined linear function.

\begin{tikzpicture}
\begin{axis}[
    xmin=-5, xmax=5,
    ymin=-5, ymax=5,
    width=4in,
    height=4in,
    xtick={-4,-3,-2,-1,0,1,2,3,4},
    ytick={-4,-3,-2,-1,0,1,2,3,4},
    tick style={line width=1.2pt},
    major tick length=6pt,
    axis x line=center,
    axis y line=center,
    x axis line style={-{Latex[length=10pt,width=8pt]}, line width=1.2pt},
    y axis line style={-{Latex[length=10pt,width=8pt]}, line width=1.2pt},
    xlabel={$x$},
    ylabel={$y$},
    every axis x label/.style={
        at={(current axis.right of origin)},
        anchor=west
    }
]

% Graph y = F(x)
\addplot[
    blue,
    line width=3pt,
    <->,
]
coordinates {
    (-4,3)
    (1,-2)
    (5,3)
};

% Points
\addplot[
    only marks,
    mark=*,
    mark size=4pt
]
coordinates {
    (1,-2)
    (0,-1)
    (-1,0)
    (2.6,0)
};

% Labels (kept clear of graph and axes)

\node at (axis cs:0,-1)
    [below left=8pt]
    {\Large $(0,-1)$};

\node at (axis cs:1,-2)
    [below=8pt]
    {\Large $(1,-2)$};

\node at (axis cs:-1,0)
    [above=8pt,left=12pt]
    {\Large $(-1,0)$};

\node at (axis cs:2.6,0)
    [above=8pt,right=10pt]
    {\Large $(2.6,0)$};

\end{axis}
\end{tikzpicture}

\begin{solution}
If $F(x) = a|x-h| + k$, then for the vertex to be at $(1,-2)$, $h  =1$ and $k = -2$ so $F(x) = a |x-1| - 2$.  Since $(0,-1)$ is on the graph, $F(0) = -1$ so $-1 = a|0-1|-2$ which means $a = 1$.  This means $F(x) = |x-1|-2$.  However, $(2.6,0)$ is also on the graph, so it should work out that $F(2.6) = 0$.  However, we find $F(2.6) = |2.6-1| - 2 = -0.4 \neq 0$.  

 \[ F(x) =  \begin{cases} 
 -x-1 &  \text{if $x \leq 1$, } \\
   \frac{5}{4} x - \frac{13}{4}  & \text{if $x \geq 1$,} \\
  \end{cases}\]
  
\end{solution}

\end{problem}

\begin{question}
In Exercises \ref{graphadvabsvalexerfirst} - \ref{graphadvabsvalexerlast}, graph the function by rewriting each function as a piecewise defined function using Definition \ref{absolutevaluepiecewise}. Find the axis intercepts of each graph, if any exist.  From the graph, determine the domain and range of each function, the maximum and minimum of each function, if they exist, and list the intervals on which the function is increasing, decreasing or constant.

\begin{problem}\label{graphadvabsvalexerfirst}
$f(x) = x + |x| - 3$  

\begin{solution}
Re-write $f(x) = x+|x| - 3$ as \\ ${\displaystyle f(x) = \left\{ \begin{array}{rcl}
-3 & \mbox{ if } & x < 0\\
     2x -3 & \mbox{ if } & x \geq 0 \\ \end{array} \right. }$ \\  $x$-intercept $\left(\frac{3}{2}, 0\right)$ \\ $y$-intercept $(0,-3)$ \\ Domain $(-\infty, \infty)$ \\ Range $[-3, \infty)$ \\ Increasing on $[0,\infty)$ \\ Constant on $(-\infty, 0]$ \\  Minimum is $-3$ at $(x,-3)$ where $x \leq 0$ \\ No  maximum  \\

\graph{x+abs(x)+3}
\end{solution}
\end{problem}

\begin{problem}
$f(x) = |x+2| - x$
\end{problem}

\begin{problem}
$f(x) = |x+2| - |x|$
\end{problem}

\begin{problem}
$g(t) = |t+ 4| + |t- 2|$
\end{problem}

\begin{problem}
$g(t) = \frac{|t + 4|}{t + 4}$
\end{problem}

\begin{problem}\label{graphadvabsvalexerlast}
$g(t) = \frac{|2 - t|}{2 - t}$
\end{problem}
  

\end{question}

\begin{problem}
With the help of your classmates, find an absolute value function whose graph is given below.

\begin{center}
\begin{tikzpicture}
\begin{axis}[
    xmin=-10, xmax=10,
    ymin=-2, ymax=6,
    width=4in,
    height=2in,
    xtick={-8,-7,-6,-5,-4,-3,-2,-1,1,2,3,4,5,6,7,8},
    ytick={1,2,3,4},
    tick style={line width=1.2pt},
    major tick length=6pt,
    axis x line=center,
    axis y line=center,
    x axis line style={-{Latex[length=10pt,width=8pt]}, line width=1.2pt},
    y axis line style={-{Latex[length=10pt,width=8pt]}, line width=1.2pt},
    xlabel={$x$},
    ylabel={$y$},
    every axis x label/.style={
        at={(current axis.right of origin)},
        anchor=west
    }
]

% Left ray
\addplot[
    blue,
    line width=3pt,
    ->,
]
coordinates {
    (-4,0)
    (-9,5)
};

% Right ray
\addplot[
    blue,
    line width=3pt,
    ->,
]
coordinates {
    (4,0)
    (9,5)
};

% Middle V-shape
\addplot[
    blue,
    line width=3pt
]
coordinates {
    (-4,0)
    (0,4)
    (4,0)
};

% Points
\addplot[
    only marks,
    mark=*,
    mark size=4pt
]
coordinates {
    (-4,0)
    (0,4)
    (4,0)
};

% Labels
\node at (axis cs:-4,0)
    [below=10pt]
    {\Large $(-4,0)$};

\node at (axis cs:0,4)
    [right=8pt]
    {\Large $(0,4)$};

\node at (axis cs:4,0)
    [below=10pt]
    {\Large $(4,0)$};

\end{axis}
\end{tikzpicture}
\end{center}
\end{problem}

\begin{question}
In Exercises \ref{solveabsvalequfirsta} - \ref{solveabsvalequlasta}, solve the equation.

\begin{problem}\label{solveabsvalequfirsta}
$|x| = 6$
\end{problem}

\begin{problem}
$|3x-1| = 10$
\end{problem}

\begin{problem}
$|4-x| = 7$
\end{problem}

\begin{problem}
$4 - |t| = 3$
\end{problem}

\begin{problem}
$2|5t+1| - 3 = 0$
\end{problem}

\begin{problem}
$|7t-1| + 2 = 0$
\end{problem}
  
\begin{problem}
$\frac{5 - |w|}{2} = 1$
\end{problem}

\begin{problem}
$\frac{2}{3} |5-2w| - \frac{1}{2} = 5$
\end{problem}

\begin{problem}
$|w| = w + 3$
\end{problem}

\begin{problem}
$|2x-1| = x+1$
\end{problem}

\begin{problem}
$4 - |x| = 2x+1$
\end{problem}

\begin{problem}\label{solveabsvalequlasta}
$|x-4| = x-5$
\end{problem}  
\end{question}

\begin{question}
Solve the equations in Exercises \ref{moreabsvalequfirsta} - \ref{moreabsvalequlasta} using  the property that if $|a| = |b|$ then  $a = \pm b$. 


\begin{problem}\label{moreabsvalequfirsta}
$|3x - 2| = |2x + 7|$ 
\end{problem}

\begin{problem}
$|3x+1| = |4x|$
\end{problem}

\begin{problem}
$|1-2x| = |x+1|$
\end{problem}

\begin{problem}
$|4-t| - |t+2| = 0$
\end{problem}

\begin{problem}
$|2-5t| = 5 |t+1|$
\end{problem}

\begin{problem}\label{moreabsvalequlasta}
$3|t-1| = 2|t+1|$
\end{problem}

\end{question}

\begin{question}
In Exercises \ref{solveinequabsfirst} - \ref{solveinequabslast}, solve the inequality.  Write your answer using interval notation. 

\begin{problem}\label{solveinequabsfirst}
$|3x - 5| \leq 4$
\end{problem}

\begin{problem}
$|7x + 2| > 10$
\end{problem}

\begin{problem}
$|2t+1| - 5 < 0$
\end{problem}

\begin{problem}
$|2-t| - 4 \geq -3$
\end{problem}

\begin{problem}
$|3w+5| + 2 < 1$
\end{problem}

\begin{problem}
$2|7-w| +4 > 1$
\end{problem}

\begin{problem}
$2 \leq |4-x| < 7$ 
\end{problem}

\begin{problem}
$1 < |2x - 9| \leq 3$
\end{problem}

\begin{problem}
$|t + 3| \geq |6t + 9|$
\end{problem}

\begin{problem}
$|t-3| - |2t+1| < 0$ 
\end{problem}

\begin{problem}
$|1-2x| \geq x + 5$
\end{problem}

\begin{problem}
$x + 5 <  |x+5|$
\end{problem}

\begin{problem}
$x \geq |x+1|$ 
\end{problem}

\begin{problem}
$|2x + 1| \leq 6-x$ 
\end{problem}

\begin{problem}
$t + |2t-3| < 2$
\end{problem}

\begin{problem}\label{solveinequabslast}
$|3-t| \geq t-5$
\end{problem}
\end{question}

\begin{problem}
Show that if $\delta$ is a real number with $\delta > 0$, the solution to $|x-a| < \delta$ is the interval:  $(a - \delta, a + \delta)$.  That is, an interval centered at $a$ with `radius' $\delta$.  
\end{problem}  

\begin{problem}\label{triangleinequalityreals}
The \href{http://en.wikipedia.org/wiki/Triangle_inequality}{\underline{Triangle Inequality}}  for real numbers states that for all real numbers $x$ and $a$, $|x+a| \leq |x| + |a|$ and, moreover, $|x+a| = |x|+|a|$ if and only if $x$ and $a$ are both positive, both negative, or one or the other is $0$.  Graph each pair of functions below on the same pair of axes and use the graphs to verify the triangle inequality in each instance.

\begin{itemize}

\item  $f(x) = |x+2|$ and $g(x) = |x|+2$.

\item  $f(x) = |x+4|$ and $g(x) = |x|+4$.

\end{itemize}

\end{problem}

\end{document}
